\section{Why a finite angular ansatz does not close}
\label{app:app-obstructions}

Section~\ref{sec:scalar-core} ruled out a monotone flat end for the
linearized height. A different tempting simplification is to keep a
finite number of angular modes in the nonlinear height. The
obstruction below explains why the inverse must accommodate all even
frequencies, even though its initial seed has only frequency two.
It is independent of the existence argument.

\begin{proposition}\label{prop:obs-finite-square}
Let $E>0$ be $C^1$ on a connected interval $I\subset(0,\infty)$.
If a real $C^2$ solution of the Darboux equation for
$E\,dr^2+r^2\,d\theta^2$ satisfies
\[
 z^2=\sum_{|j|\le J}b_j(r)e^{ij\theta},\qquad
 b_{-j}=\overline{b_j},\qquad b_j\in C^2(I),
\]
and is radial on a nonempty open radial interval, then it is radial
throughout $I$. This includes every finite Fourier height, and every
smooth square root of a finite trigonometric polynomial.
\end{proposition}

The issue is continuation through a zero of a Fourier coefficient.
A nonzero coefficient has a complex phase, so a real-valued
uniqueness argument would not suffice. We give the continuation
argument rather than assume that the resulting singular ordinary
differential equation has ordinary uniqueness.

\begin{proof}
Set $v=z^2$ and $K=E'/(2rE^2)$. The Hessian identity
$\nabla^2v=2z\nabla^2z+2dz\otimes dz$, together with the
two-dimensional rank-one determinant identity, gives
\begin{equation}\label{eq:obstruction-polynomial}
 2v\det_g\nabla^2v-\operatorname{cof}_g(\nabla^2v)[dv,dv]
       =8Kv^2-2Kv|dv|_g^2.
\end{equation}
To check it without dividing by $z$, expand the two terms on the
left: the terms containing
$z^3\operatorname{cof}_g(\nabla^2z)[dz,dz]$ cancel, leaving
$8z^4\det_g\nabla^2z$. Substitution of the Darboux equation
gives the right side. Thus the identity also holds where $z=0$.

Let $b=b_J$ for the highest positive frequency $J$, and put
$a=EJ^2/r$, $d=E'J^2/(2r)$. The coefficient of frequency $3J$
in \eqref{eq:obstruction-polynomial} contains only three copies
of $b$. It is
\begin{equation}\label{eq:obstruction-highest-mode}
\begin{split}
0={}&(2bb'-ab^2)b''-(b')^3+ab(b')^2
 -\frac{2a}{r}b^2b'+\frac{2a}{r^2}b^3\\
 &\quad+db^2b'-\frac{2d}{r}b^3.
\end{split}
\end{equation}
For a direct coefficient check, write the three Hessian entries of
$v$ as
$A=v_{rr}-E'v_r/(2E)$,
$B=v_{r\theta}-v_\theta/r$,
$C=v_{\theta\theta}+rv_r/E$.
After multiplication by $Er^2$, the left side is
$2v(AC-B^2)-(Cv_r^2-2Bv_rv_\theta+Av_\theta^2)$.
Replace the jets by $b,b',b'',iJb,iJb',-J^2b$ and collect
terms; the quadratic term on the right has no frequency $3J$.

On a component where $b\ne0$, take a continuous complex square root
$y^2=b$. Substitution reduces \eqref{eq:obstruction-highest-mode} to
\begin{equation}\label{eq:obstruction-square-root}
 (4y'-ay)y''+a(y'-y/r)^2+d(yy'-y^2/r)=0.
\end{equation}
Fix a compact subinterval and $K_1>1+\sup|a|$. The quantity
\[
 U=\max(K_1|y|,|y'|)
   =\max\left(K_1\sqrt{|b|},
                    \frac{|b'|}{2\sqrt{|b|}}\right)
\]
satisfies $|U'|\le CU$ almost everywhere. Indeed, where the second
term is larger, $|4y'-ay|\ge3|y'|$, and
\eqref{eq:obstruction-square-root} gives $|y''|\le C|y'|$.
Where the first is larger, its derivative has absolute value at
most $K_1|y'|\le K_1U$. The maximum is locally Lipschitz,
so the same estimate holds through its switching points.
It bounds $U$ both above and away from zero at every finite
endpoint of a nonzero component.

At a zero of $b$, \eqref{eq:obstruction-highest-mode} gives
$b'=0$. At a finite endpoint of a nonzero component, $y\to0$,
so $U=|y'|$ nearby. The estimate for $y''$ yields a limit
$y'\to\lambda\ne0$, and consequently
\begin{equation}\label{eq:obstruction-quadratic-zero}
 b''=2(y')^2+2yy''\longrightarrow2\lambda^2\ne0.
\end{equation}
Thus such an endpoint is an isolated quadratic zero. Across it
$U$ has the continuous value $\sqrt{|b''|/2}$ and the same
differential bound.

Consider the open set
$\mathcal O=\{b\ne0\}\cup\{b=0, b''\ne0\}$.
If a component of $\mathcal O$ had an interior boundary point,
then $b=b''=0$ there. An infinite accumulation of quadratic zeros
would contradict the positive lower bound for $U$ and continuity
of $b''$; a last nonzero interval would contradict
\eqref{eq:obstruction-quadratic-zero}. Hence
$\mathcal O$ is either empty or all of $I$. Since $b$ vanishes
on the given open radial interval, it must vanish everywhere.

Remove this highest coefficient and repeat down the finite
expansion. The resulting $z^2$ is radial. At each radius a real
continuous function on the connected circle whose square is
constant is itself constant. Therefore $z$ is radial.
\end{proof}

This result does not say that a single angular mode is useless.
It is exactly the right seed for the linearized construction.
It says that the nonlinear completion cannot be confined to a
finite angular ansatz once it is required to meet a radial collar.
